\section{Bibliographic comparisons}\label{app:literature}

These notes record the scope of related constructions and restrictions.
They complement the attributions attached directly to the results above.


\subsection{A rationality filter}

A fully magic square of nine distinct squares \emph{does exist} over a
number field of degree four. Bremner displayed the underlying magic square
with eight entries square in $\Q(\sqrt3)$ and centre $532$, which is not a
square in that field~\cite{BRE01}; adjoining $\sqrt{133}$ makes the centre
$532=(2\sqrt{133})^{2}$. Thus the degree-four completion is the elementary
field extension of Bremner's displayed object, not an additional
construction attributed to him. The companion repository replays it as an
exact verified witness. Bremner~\cite{BRE99},
\S3, separately constructs an example over a field of degree $27$. The
consequence deserves emphasis: there is no
combinatorial or shape-only obstruction to the rational magic-square problem. Thus the obstruction to a rational solution must use arithmetic
information that distinguishes $\Q$ from the displayed quartic field.
This filter shapes the whole paper: the positive constructions of
Section~\ref{sec:families} show how close the rational story comes, and the
support laws of Sections~\ref{sec:threeprime} and~\ref{sec:supportlaw} are
genuinely rational-arithmetic statements, invisible over a larger field.

\subsection{The entry side: what is known about a hypothetical solution}

A hypothetical rational solution is constrained through the prime factors of
its \emph{entries}; throughout this paragraph the square is taken in the
primitive normalization (unit gcd of the nine entries).
Rabern~\cite{RAB03} proved --- Weisenberg's restatement in~\cite{WEI23},
\S2.1, matches the primary on all six results --- that every prime factor of
an entry is odd; that a prime dividing the centre is $1 \bmod 4$; that a
prime $p\equiv3,5 \pmod 8$ dividing a non-centre entry also divides the
centre and the third entry of that line; that no $p\equiv3\pmod 8$ divides
any entry; that no $p\equiv5\pmod8$ divides a middle-side entry; and that a
$p\equiv3\pmod4$ dividing a corner divides the two non-adjacent middle-side
entries. Woll~\cite{WOL18} proved that a prime $p$ dividing only the centre
forces $p$ to possess three consecutive nonzero quadratic residues, and
counted the resulting residue classes; Labruna~\cite{LAB18} studied the
order-three analogue over finite fields, where in characteristic greater
than three a nontrivial magic square of squares attains exactly $3$, $5$,
$7$ or $9$ distinct entries, each degree realized over infinitely many
prime fields. Weisenberg~\cite{WEI23} added the
order bookkeeping (a prime dividing exactly three entries attains its
minimum order at least twice), the restriction $p\equiv1\pmod{24}$ for a
prime dividing exactly two entries, consecutive-residue patterns for
single-entry primes, a complete seven-arrangement list of the ways a prime
factor can appear, and a proposal --- not executed there --- to organize
searches by the Gaussian prime signature of the centre. A recent preprint
of Hill~\cite{HILL26} announces a nonexistence proof; it is unrefereed,
and nothing in this paper depends on it.
Coumbe~\cite{COU24} develops a Pell-transform approach and claims that the
smallest entry cannot be unity or a prime square. We do not import that
exclusion: the unrestricted difference-set statement used in its proof has
the literal instance $120+120=240$, coming from the two primitive
progressions $(7^2,13^2,17^2)$ and $(7^2,17^2,23^2)$. This does not refute
a repaired distinct-summand statement or the final exclusion itself, but it
prevents the published proof from being used here without repair.

All of these statements constrain \emph{entries of a completed square}.
The present paper works on a different quantity --- the common difference
of the three square progressions forced by the completion problem --- and
Section~\ref{sec:supportlaw} depends on none of the entry-side results;
the two layers meet only in the classical toolbox of quadratic residues
and sums of two squares.



\subsection{Provenance of the centre-true eight-square arrays}

A related catalogued family must not be confused with the $(8,7)$ metric:
the \emph{centre-true} arrays with a nonsquare centre $g$, eight distinct
square entries around it, magic sum $3g$, and \textbf{six} magic lines
(three rows, middle column, both diagonals; the outer columns fail). Seven
such arrays with centres
\[
157441,\;411625,\;786361,\;1199809,\;2033641,\;12138289,\;19953649
\]
appear in Brown's note on automedian triangles and magic
squares~\cite{BRO} which lists the six roots of the top and bottom rows for each centre.
The square middle-pair completion gives the six-line arrangement below.
We include it to distinguish this catalogued construction from our
$(8,7)$ metric, without treating the centres or their outer rows as new. For example, at $g=157441$:
\[
\begin{pmatrix}
421^2 & 49^2 & 541^2\\
481^2 & g & 289^2\\
149^2 & 559^2 & 371^2
\end{pmatrix},
\]
whose six stated lines each sum to $3g=472323$. Brown also states that a
Parker centre root needs more than two distinct primes $1\bmod4$; that
statement is non-refereed, its proof scope has not been established here,
and it is not an input to any of our proofs.



\subsection{Progressions in elliptic coordinate sets}

Harrison--Mudgal--Schmidt~\cite{HMS26} give uniform rank-dependent bounds
for progressions inside one elliptic coordinate set. This is complementary
to the three-coordinate interaction studied here and does not decide its
positive distinct-entry locus. The fixed-squareclass finiteness proof in
Section~\ref{sec:finiteness} instead applies the explicit classification
of elliptic curves in a linear $x$-coordinate relation due to
Caro--Garc\'ia-Fritz~\cite{CGF25}.
